
\documentclass{article}
\usepackage{colortbl}
\usepackage{makecell}
\usepackage{multirow}
\usepackage{supertabular}

\begin{document}

\newcounter{utterance}

\centering \large Interaction Transcript for game `cladder', experiment `full\_v1.5\_default', episode 905 with Qwen3.6{-}27B.
\vspace{24pt}

{ \footnotesize  \setcounter{utterance}{1}
\setlength{\tabcolsep}{0pt}
\begin{supertabular}{c@{$\;$}|p{.15\linewidth}@{}p{.15\linewidth}p{.15\linewidth}p{.15\linewidth}p{.15\linewidth}p{.15\linewidth}}
    \# & \multicolumn{2}{c}{Player} && \multicolumn{2}{c}{Game Master} \\
    \hline

    \theutterance \stepcounter{utterance}  
    & & & \multicolumn{4}{p{0.6\linewidth}}{
        \cellcolor[rgb]{0.9,0.9,0.9}{
            \makecell[{{p{\linewidth}}}]{
                \texttt{\tiny{[P1$\langle$GM]}}
                \texttt{You are an expert in causal inference. The following question is not a typical commonsense query, but rather a meticulously designed question created by a professor specializing in causal inference, intended to assess the students' mastery of the course content.} \\
\texttt{Imagine a self{-}contained, hypothetical world with only the following conditions, and without any unmentioned factors or causal relationships: Smoking has a direct effect on effort and college admission. Effort has a direct effect on college admission. For nonsmokers, the probability of college admission is 76\%. For smokers, the probability of college admission is 22\%. For smokers, would it be less likely to see college admission if the student had been a nonsmoker?} \\
\texttt{Start your answer with 'Yes' or 'No', followed by additional reasoning or evidence to support your explanation.} \\
            }
        }
    }
     \\ \\

    \theutterance \stepcounter{utterance}  
    & \multicolumn{4}{p{0.6\linewidth}}{
        \cellcolor[rgb]{0.561,0.737,0.561}{
            \makecell[{{p{\linewidth}}}]{
                \texttt{\tiny{[P1$\rangle$GM]}}
                \texttt{No} \\
\\ 
\texttt{**Reasoning:**} \\
\\ 
\texttt{To answer this question, we must distinguish between the **associational** probability (what we observe) and the **counterfactual** probability (what would have happened under a different intervention).} \\
\\ 
\texttt{1.  **Analyze the Causal Structure:**} \\
\texttt{    The problem states:} \\
\texttt{    *   Smoking (\$S\$) \$\textbackslash{}rightarrow\$ Effort (\$E\$)} \\
\texttt{    *   Smoking (\$S\$) \$\textbackslash{}rightarrow\$ College Admission (\$A\$)} \\
\texttt{    *   Effort (\$E\$) \$\textbackslash{}rightarrow\$ College Admission (\$A\$)} \\
\\ 
\texttt{    This implies that Smoking affects Admission through two pathways:} \\
\texttt{    *   A direct effect: \$S \textbackslash{}rightarrow A\$} \\
\texttt{    *   An indirect effect via Effort: \$S \textbackslash{}rightarrow E \textbackslash{}rightarrow A\$} \\
\\ 
\texttt{2.  **Analyze the Observed Data:**} \\
\texttt{    *   \$P(A=1 | S=0) = 0.76\$ (Probability of admission for nonsmokers)} \\
\texttt{    *   \$P(A=1 | S=1) = 0.22\$ (Probability of admission for smokers)} \\
\\ 
\texttt{    The observed difference is \$0.22 {-} 0.76 = {-}0.54\$. This means smokers are observed to be 54 percentage points *less* likely to be admitted than nonsmokers.} \\
\\ 
\texttt{3.  **Analyze the Counterfactual Question:**} \\
\texttt{    The question asks} \\
            }
        }
    }
    & & \\ \\

    \theutterance \stepcounter{utterance}  
    & & & \multicolumn{2}{p{0.3\linewidth}}{
        \cellcolor[rgb]{0.9,0.9,0.9}{
            \makecell[{{p{\linewidth}}}]{
                \texttt{\tiny{[GM$|$GM]}}
                \texttt{No} \\
            }
        }
    }
    & & \\ \\

    \theutterance \stepcounter{utterance}  
    & & & \multicolumn{2}{p{0.3\linewidth}}{
        \cellcolor[rgb]{0.9,0.9,0.9}{
            \makecell[{{p{\linewidth}}}]{
                \texttt{\tiny{[GM$|$GM]}}
                \texttt{no} \\
            }
        }
    }
    & & \\ \\

    \theutterance \stepcounter{utterance}  
    & & & \multicolumn{2}{p{0.3\linewidth}}{
        \cellcolor[rgb]{0.9,0.9,0.9}{
            \makecell[{{p{\linewidth}}}]{
                \texttt{\tiny{[GM$|$GM]}}
                \texttt{game\_result = WIN} \\
            }
        }
    }
    & & \\ \\

\end{supertabular}
}

\end{document}
